\documentclass{amsart}
% For dealing with references we use the comment environment
\usepackage{verbatim}
\newenvironment{reference}{\comment}{\endcomment}
%\newenvironment{reference}{}{}
\newenvironment{slogan}{\comment}{\endcomment}
\newenvironment{history}{\comment}{\endcomment}

% For commutative diagrams we use Xy-pic
\usepackage[all]{xy}

% We use 2cell for 2-commutative diagrams.
\xyoption{2cell}
\UseAllTwocells

% We use multicol for the list of chapters between chapters
\usepackage{multicol}

% This is generally recommended for better output
\usepackage{lmodern}
\usepackage[T1]{fontenc}

% For cross-file-references

% Package for hypertext links:
\usepackage{hyperref}

% For any local file, say "hello.tex" you want to link to please

% Theorem environments.
%
\theoremstyle{plain}
\newtheorem{theorem}[subsection]{Theorem}
\newtheorem{proposition}[subsection]{Proposition}
\newtheorem{lemma}[subsection]{Lemma}

\theoremstyle{definition}
\newtheorem{definition}[subsection]{Definition}
\newtheorem{example}[subsection]{Example}
\newtheorem{exercise}[subsection]{Exercise}
\newtheorem{situation}[subsection]{Situation}

\theoremstyle{remark}
\newtheorem{remark}[subsection]{Remark}
\newtheorem{remarks}[subsection]{Remarks}

\numberwithin{equation}{subsection}

% Macros
%
\def\lim{\mathop{\mathrm{lim}}\nolimits}
\def\colim{\mathop{\mathrm{colim}}\nolimits}
\def\Spec{\mathop{\mathrm{Spec}}}
\def\Hom{\mathop{\mathrm{Hom}}\nolimits}
\def\Ext{\mathop{\mathrm{Ext}}\nolimits}
\def\SheafHom{\mathop{\mathcal{H}\!\mathit{om}}\nolimits}
\def\SheafExt{\mathop{\mathcal{E}\!\mathit{xt}}\nolimits}
\def\Sch{\mathit{Sch}}
\def\Mor{\mathop{\mathrm{Mor}}\nolimits}
\def\Ob{\mathop{\mathrm{Ob}}\nolimits}
\def\Sh{\mathop{\mathit{Sh}}\nolimits}
\def\NL{\mathop{N\!L}\nolimits}
\def\CH{\mathop{\mathrm{CH}}\nolimits}
\def\proetale{{pro\text{-}\acute{e}tale}}
\def\etale{{\acute{e}tale}}
\def\QCoh{\mathit{QCoh}}
\def\Ker{\mathop{\mathrm{Ker}}}
\def\Im{\mathop{\mathrm{Im}}}
\def\Coker{\mathop{\mathrm{Coker}}}
\def\Coim{\mathop{\mathrm{Coim}}}

% Boxtimes
%
\DeclareMathSymbol{\boxtimes}{\mathbin}{AMSa}{"02}

%
% Macros for moduli stacks/spaces
%
\def\QCohstack{\mathcal{QC}\!\mathit{oh}}
\def\Cohstack{\mathcal{C}\!\mathit{oh}}
\def\Spacesstack{\mathcal{S}\!\mathit{paces}}
\def\Quotfunctor{\mathrm{Quot}}
\def\Hilbfunctor{\mathrm{Hilb}}
\def\Curvesstack{\mathcal{C}\!\mathit{urves}}
\def\Polarizedstack{\mathcal{P}\!\mathit{olarized}}
\def\Complexesstack{\mathcal{C}\!\mathit{omplexes}}
% \Pic is the operator that assigns to X its picard group, usage \Pic(X)
% \Picardstack_{X/B} denotes the Picard stack of X over B
% \Picardfunctor_{X/B} denotes the Picard functor of X over B
\def\Pic{\mathop{\mathrm{Pic}}\nolimits}
\def\Picardstack{\mathcal{P}\!\mathit{ic}}
\def\Picardfunctor{\mathrm{Pic}}
\def\Deformationcategory{\mathcal{D}\!\mathit{ef}}

\setlength{\emergencystretch}{3em}
\begin{document}
\title[Stacks Project, Tag 05D6]{1311 --- Countably generated Mittag-Leffler modules over henselian local rings split into finitely presented summands}
\maketitle
\noindent Copyright (C) 2005--2025 Johan de Jong. Stacks Project authors. Source: \href{https://stacks.math.columbia.edu/tag/05D6}{Tag 05D6}.

\noindent Editorial difficulty note (not author text). Difficulty H2: The proof first produces finite-presentation-factorizing endomorphisms fixing prescribed elements, then turns an annihilating characteristic polynomial into a henselian spectral factorization and an actual direct-sum decomposition. Iterating these polynomially constructed summands yields an infinite decomposition; this combines infinite-module approximation and henselian algebra beyond qualification arguments..

\noindent Editorial provenance note (not author text). History: excerpt from revision \texttt{a0444\allowbreak 6e57e\allowbreak c1fbc\allowbreak 252a8\allowbreak 71afc\allowbreak ec775\allowbreak 2fb28\allowbreak 07b14}, algebra.tex, lines 42692--42757. Reference presentation was adapted; original-author context and core support blocks were retained or added. The mathematical wording is unchanged. Licensed under GNU FDL 1.2 or later; no invariant sections or cover texts. See ../licenses/COPYING. Context blocks reproduce author definitions and supporting results; they are not additional counted problems.

\begin{definition}
\label{definition-mittag-leffler-module}
Let $M$ be an $R$-module.  We say that $M$ is {\it Mittag-Leffler} if the
equivalent conditions of
Proposition \ref{proposition-ML-characterization}
hold.
\end{definition}

\begin{definition}
\label{definition-domination}
Let $f: M \to N$ and $g: M \to M'$ be maps of $R$-modules.
Then we say $g$ {\it dominates} $f$ if for any $R$-module $Q$, we have $\Ker(f
\otimes_R \text{id}_Q) \subset \Ker(g \otimes_R \text{id}_Q)$.
\end{definition}

\begin{definition}
\label{definition-universally-injective}
Let $f: M \to N$ be a map of $R$-modules.  Then $f$ is called
{\it universally injective} if for every $R$-module $Q$, the map $f
\otimes_R \text{id}_Q: M \otimes_R Q \to N \otimes_R Q$
is injective.  A sequence $0 \to M_1 \to M_2 \to M_3
\to 0$ of $R$-modules is called {\it universally exact} if it is exact
and $M_1 \to M_2$ is universally injective.
\end{definition}

\begin{lemma}
\label{lemma-ML-countable}
Let $R$ be a ring.
Let $M$ be an $R$-module.
Assume $M$ is Mittag-Leffler and countably generated.
For any $R$-module map $f : P \to M$ with $P$ finitely generated there
exists an endomorphism $\alpha : M \to M$ such that
\begin{enumerate}
\item $\alpha : M \to M$ factors through a finitely presented $R$-module, and
\item $\alpha \circ f = f$.
\end{enumerate}
\end{lemma}

\begin{proof}
Write $M = \colim_{i \in I} M_i$ as a directed colimit of finitely
presented $R$-modules with $I$ countable, see
Lemma \ref{lemma-ML-countable-colimit}.
The transition maps are denoted $f_{ij}$ and we use $f_i : M_i \to M$
to denote the canonical maps into $M$. Set $N = \prod_{s \in I} M_s$. Denote
$$
M_i^* = \Hom_R(M_i, N) = \prod\nolimits_{s \in I} \Hom_R(M_i, M_s)
$$
so that $(M_i^*)$ is an inverse system of $R$-modules over $I$.
Note that $\Hom_R(M, N) = \lim M_i^*$.
As $M$ is Mittag-Leffler, we find for every
$i \in I$ an index $k(i) \geq i$ such that
$$
E_i := \bigcap\nolimits_{i' \geq i} \Im(M_{i'}^* \to M_i^*)
=
\Im(M_{k(i)}^* \to M_i^*)
$$
Choose and fix $j \in I$ such that $\Im(P \to M) \subset \Im(M_j \to M)$.
This is possible as $P$ is finitely generated. Set $k = k(j)$.
Let
$x = (0, \ldots, 0, \text{id}_{M_k}, 0, \ldots, 0) \in M_k^*$ and
note that this maps to $y = (0, \ldots, 0, f_{jk}, 0, \ldots, 0) \in M_j^*$.
By our choice of $k$ we see that $y \in E_j$. By
Example \href{https://stacks.math.columbia.edu/tag/0596}{example ML surjective maps (Tag 0596)}
the transition maps $E_i \to E_j$ are surjective for each $i \geq j$
and $\lim E_i = \lim M_i^* = \Hom_R(M, N)$. Hence
Lemma \href{https://stacks.math.columbia.edu/tag/0597}{lemma ML limit nonempty (Tag 0597)}
guarantees there exists an element $z \in \Hom_R(M, N)$
which maps to $y$ in $E_j \subset M_j^*$. Let $z_k$ be the $k$th component
of $z$. Then $z_k : M \to M_k$ is a homomorphism such that
$$
\xymatrix{
M \ar[r]_{z_k} & M_k \\
M_j \ar[ru]_{f_{jk}} \ar[u]^{f_j}
}
$$
commutes. Let $\alpha : M \to M$ be the composition
$f_k \circ z_k : M \to M_k \to M$.
Then $\alpha$ factors through a finitely presented module by construction and
$\alpha \circ f_j = f_j$. Since the image of $f$ is contained in the image of
$f_j$ this also implies that $\alpha \circ f = f$.
\end{proof}


\begin{lemma}
\label{lemma-ML-countable-colimit}
Let $M$ be an $R$-module.  Write $M = \colim_{i \in I} M_i$ where $(M_i,
f_{ij})$ is a directed system of finitely presented $R$-modules.  If $M$ is
Mittag-Leffler and countably generated, then there is a directed countable
subset $I' \subset I$ such that $M \cong \colim_{i \in I'} M_i$.
\end{lemma}

\begin{proof}
Let $x_1, x_2, \ldots$ be a countable set of generators for $M$.  For each $x_n$
choose $i \in I$ such that $x_n$ is in the image of the canonical map $f_i:
M_i \to M$; let $I'_{0} \subset I$ be the set of all these $i$.  Now
since $M$ is Mittag-Leffler, for each $i \in I'_{0}$ we can choose $j \in I$
such that $j \geq i$ and $f_{ij}: M_i \to M_j$ factors through
$f_{ik}: M_i \to M_k$ for all $k \geq i$  (condition (3) of Proposition
\ref{proposition-ML-characterization}); let $I'_1$ be the union of $I'_0$ with
all of these $j$.  Since $I'_1$ is a countable, we can enlarge it to a
countable directed set $I'_{2} \subset I$.  Now we can apply the same procedure
to $I'_{2}$ as we did to $I'_{0}$ to get a new countable set $I'_{3} \subset
I$.  Then we enlarge $I'_{3}$ to a countable directed set $I'_{4}$.  Continuing
in this way---adding in a $j$ as in Proposition
\ref{proposition-ML-characterization} (3) for each $ i \in I'_{\ell}$ if $\ell$
is odd and enlarging $I'_{\ell}$ to a directed set if $\ell$ is even---we get a
sequence of subsets $I'_{\ell} \subset I$ for $\ell \geq 0$.  The union $I' =
\bigcup I'_{\ell}$ satisfies:
\begin{enumerate}
\item $I'$ is countable and directed;
\item each $x_n$ is in the image of $f_i: M_i \to M$ for some $i
\in I'$;
\item if $i \in I'$, then there is $j \in I'$ such that $j \geq i$ and $f_{ij}:
M_i \to M_j$ factors through $f_{ik}: M_i \to M_k$ for all
$k \in I$ with $k \geq i$.  In particular $\Ker(f_{ik}) \subset \Ker(f_{ij})$
for $k \geq i$.
\end{enumerate}
We claim that the canonical map $\colim_{i \in I'} M_i \to
\colim_{i
\in I} M_i = M$ is an isomorphism.  By (2) it is surjective.  For injectivity,
suppose $x \in \colim_{i \in I'} M_i$ maps to $0$ in $\colim_{i \in
I} M_i$.
Representing $x$ by an element $\tilde{x} \in M_i$ for some $i \in I'$, this
means that $f_{ik}(\tilde{x}) = 0$ for some $k \in I, k \geq i$.  But then by
(3) there is $j \in I', j \geq i,$ such that $f_{ij}(\tilde{x}) = 0$.  Hence $x
= 0$ in $\colim_{i \in I'} M_i$.
\end{proof}


\begin{proposition}
\label{proposition-ML-characterization}
Let $M$ be an $R$-module.  Let $(M_i, f_{ij})$ be a directed system of finitely
presented $R$-modules, indexed by $I$, such that $M = \colim M_i$.  Let
$f_i:
M_i \to M$ be the canonical map.  The following are equivalent:
\begin{enumerate}
\item For every finitely presented $R$-module $P$ and module map $f: P
\to M$, there exists a finitely presented $R$-module $Q$ and a module
map $g: P \to Q$ such that $g$ and $f$ dominate each other, i.e.,
$\Ker(f \otimes_R \text{id}_N) = \Ker(g \otimes_R \text{id}_N)$
for every $R$-module $N$.
\item For each $i \in I$, there exists $j \geq i$ such that $f_{ij}: M_i
\to M_j$ dominates $f_i: M_i \to M$.
\item For each $i \in I$, there exists $j \geq i$ such that $f_{ij}: M_i
\to M_j$ factors through $f_{ik}: M_i \to M_k$ for all $k \geq
i$.
\item For every $R$-module $N$, the inverse system
$(\Hom_R(M_i, N), \Hom_R(f_{ij}, N))$ is Mittag-Leffler.
\item For $N = \prod_{s \in I} M_s$, the inverse system
$(\Hom_R(M_i, N), \Hom_R(f_{ij}, N))$ is Mittag-Leffler.
\end{enumerate}
\end{proposition}

\begin{proof}
First we prove the equivalence of (1) and (2).  Suppose (1) holds and let $i
\in I$. Corresponding to the map $f_i: M_i \to M$, we can choose $g:
M_i \to Q$ as in (1).  Since $M_i$ and $Q$ are of finite presentation,
so is $\Coker(g)$.  Then by Lemma \ref{lemma-domination}, $f_i : M_i
\to M$ factors through $g: M_i \to Q$, say $f_i = h \circ g$
for some $h: Q \to M$.  Then since $Q$ is finitely presented, $h$
factors through $M_j \to M$ for some $j \geq i$, say $h = f_j \circ h'$
for some $h': Q \to M_j$.  In total we have a commutative diagram
$$
\xymatrix{
  &  M  & \\
M_i \ar[dr]_g \ar[ur]^{f_i} \ar[rr]^{f_{ij}} &
& M_j \ar[ul]_{f_j} \\
  & Q  \ar[ur]_{h'} &
}
$$
Thus $f_{ij}$ dominates $g$.  But $g$ dominates $f_i$, so $f_{ij}$ dominates
$f_i$.

\medskip\noindent
Conversely, suppose (2) holds.  Let $P$ be of finite presentation and $f: P
\to M$ a module map.  Then $f$ factors through $f_i: M_i \to M$
for some $i \in I$, say $f = f_i \circ g'$ for some $g': P \to M_i$.
Choose by (2) a $j \geq i$ such that $f_{ij}$ dominates $f_i$.  We have a
commutative diagram
$$
\xymatrix{
P \ar[d]_{g'} \ar[r]^{f}           & M  \\
M_i \ar[ur]^{f_i} \ar[r]_{f_{ij}} & M_j \ar[u]_{f_j}
}
$$
From the diagram and the fact that $f_{ij}$ dominates $f_i$, we find that $f$
and $f_{ij} \circ g'$ dominate each other.  Hence taking $g = f_{ij} \circ g' :
P \to M_j$ works.

\medskip\noindent
Next we prove (2) is equivalent to (3).  Let $i \in I$.  It is always true that
$f_i$ dominates $f_{ik}$ for $k \geq i$, since $f_i$ factors through
$f_{ik}$.  If (2) holds, choose $j \geq i$ such that $f_{ij}$ dominates
$f_i$.  Then since domination is a transitive relation, $f_{ij}$ dominates
$f_{ik}$ for $k \geq i$. All $M_i$ are of finite presentation, so
$\Coker(f_{ik})$ is of finite presentation for $k \geq i$.  By Lemma
\ref{lemma-domination}, $f_{ij}$ factors through $f_{ik}$ for all $k \geq i$.
Thus (2) implies (3).  On the other hand, if (3) holds then for any $R$-module
$N$, $f_{ij} \otimes_R \text{id}_N$ factors through $f_{ik}
\otimes_R \text{id}_N$ for $k \geq i$.  So $\Ker(f_{ik} \otimes_R
\text{id}_N) \subset \Ker(f_{ij} \otimes_R \text{id}_N)$ for $k
\geq i$.  But $\Ker(f_i \otimes_R \text{id}_N: M_i \otimes_R N
\to M \otimes_R N)$ is the union of $\Ker(f_{ik} \otimes_R
\text{id}_N)$ for $k \geq i$.  Thus $\Ker(f_i \otimes_R
\text{id}_N) \subset \Ker(f_{ij} \otimes_R \text{id}_N)$ for
any $R$-module $N$, which by definition means $f_{ij}$ dominates $f_i$.

\medskip\noindent
It is trivial that (3) implies (4) implies (5).  We show (5) implies (3).  Let
$N = \prod_{s \in I} M_s$. If (5) holds, then given $i \in I$ choose $j \geq i$
such that
$$
\Im( \Hom(M_j, N) \to  \Hom(M_i, N)) =
\Im( \Hom(M_k, N) \to  \Hom(M_i, N))
$$
for all $k \geq j$.  Passing the product over $s \in I$ outside of the
$\Hom$'s
and looking at the maps on each component of the product, this says
$$
\Im( \Hom(M_j, M_s) \to  \Hom(M_i, M_s)) =
\Im( \Hom(M_k, M_s) \to   \Hom(M_i, M_s))
$$
for all $k \geq j$ and $s \in I$.  Taking $s = j$ we have
$$
\Im( \Hom(M_j, M_j) \to  \Hom(M_i, M_j)) =
\Im( \Hom(M_k, M_j) \to   \Hom(M_i, M_j))
$$
for all $k \geq j$.  Since $f_{ij}$ is the image of
$\text{id} \in  \Hom(M_j, M_j)$ under
$\Hom(M_j, M_j) \to  \Hom(M_i, M_j)$,
this shows  that for any $k \geq j$ there is $h \in \Hom(M_k, M_j)$
such that $f_{ij} = h \circ f_{ik}$.  If $j \geq k$ then we can take
$h = f_{kj}$. Hence (3) holds.
\end{proof}


\begin{lemma}
\label{lemma-domination}
Let $f: M \to N$ and $g: M \to M'$ be maps of $R$-modules.
Suppose $\Coker(f)$ is of finite presentation.  Then $g$ dominates $f$ if
and
only if $g$ factors through $f$, i.e.\ there exists a module map $h: N
\to M'$ such that $g = h \circ f$.
\end{lemma}

\begin{proof}
Consider the pushout of $f$ and $g$ as in the statement of
Lemma \ref{lemma-domination-universally-injective}.
From the construction of the pushout it follows that
$\Coker(f') = \Coker(f)$, so $\Coker(f')$ is of finite
presentation.  Then by
Lemma \href{https://stacks.math.columbia.edu/tag/058L}{lemma universally exact split (Tag 058L)}, $f'$ is universally injective if and
only if
$$
0 \to M' \xrightarrow{f'} N' \to \Coker(f') \to 0
$$
splits. This is the case if and only if there is a map $h' : N' \to M'$
such that $h' \circ f' = \text{id}_{M'}$.  From the universal
property of the pushout, the existence of such an $h'$ is equivalent to $g$
factoring through $f$.
\end{proof}


\begin{lemma}
\label{lemma-domination-universally-injective}
Let $f : M \to N$ and $g : M \to M'$ be maps of $R$-modules.
Consider the pushout of $f$ and $g$,
$$
\xymatrix{
M  \ar[r]_f \ar[d]_g & N \ar[d]^{g'} \\
M' \ar[r]^{f'} & N'
}
$$
Then $g$ dominates $f$ if and only if $f'$ is universally injective.
\end{lemma}

\begin{proof}
Recall that $N'$ is $M' \oplus N$ modulo the submodule consisting of elements
$(g(x), -f(x))$ for $x \in M$.
From the construction of $N'$ we have a short exact sequence
$$
0 \to \Ker(f) \cap \Ker(g) \to \Ker(f) \to \Ker(f')
\to 0.
$$
Since tensoring commutes with taking pushouts, we have such a short exact
sequence
$$
0 \to \Ker(f \otimes \text{id}_Q ) \cap \Ker(g \otimes
\text{id}_Q) \to \Ker(f \otimes \text{id}_Q)
\to \Ker(f' \otimes \text{id}_Q) \to 0
$$
for every $R$-module $Q$.  So $f'$ is universally injective if and only if
$\Ker(f \otimes \text{id}_Q ) \subset \Ker(g \otimes
\text{id}_Q)$ for every $Q$, if and only if $g$ dominates $f$.
\end{proof}


\begin{lemma}
\label{lemma-domination-fp}
Let $f: M \to N$ and $g: M \to M'$ be maps of $R$-modules.
Then $g$ dominates $f$ if and only if for any finitely presented $R$-module
$Q$, we have $\Ker(f \otimes_R \text{id}_Q) \subset \Ker(g \otimes_R
\text{id}_Q)$.
\end{lemma}

\begin{proof}
Suppose $\Ker(f \otimes_R \text{id}_Q) \subset \Ker(g \otimes_R
\text{id}_Q)$ for all finitely presented modules $Q$.  If $Q$ is an
arbitrary module, write $Q = \colim_{i \in I} Q_i$ as a colimit of a
directed
system of finitely presented modules $Q_i$.  Then $\Ker(f \otimes_R
\text{id}_{Q_i}) \subset \Ker(g \otimes_R \text{id}_{Q_i})$ for
all $i$.  Since taking directed colimits is exact and commutes with tensor
product, it follows that $\Ker(f \otimes_R \text{id}_Q) \subset \Ker(g
\otimes_R \text{id}_Q)$.
\end{proof}


\begin{lemma}
\label{lemma-direct-sum-ML}
If $M = \bigoplus_{i \in I} M_i$ is a direct sum of $R$-modules, then $M$ is
Mittag-Leffler if and only if each $M_i$ is Mittag-Leffler.
\end{lemma}

\begin{proof}
The ``only if'' direction follows from Lemma \href{https://stacks.math.columbia.edu/tag/059N}{lemma pure submodule ML (Tag 059N)} (1)
and the fact that a split short exact sequence is universally exact. The
converse follows from Lemma \href{https://stacks.math.columbia.edu/tag/0AS7}{lemma colimit universally injective ML (Tag 0AS7)}
but we can also argue it directly as follows. First note that if $I$ is
finite then this follows from Lemma
\href{https://stacks.math.columbia.edu/tag/059N}{lemma pure submodule ML (Tag 059N)} (2).  For general $I$, if all $M_i$ are
Mittag-Leffler then we prove the same of $M$ by verifying condition (1) of
Proposition \ref{proposition-ML-characterization}.
Let $f: P \to M$ be a map from a finitely presented module $P$.
Then $f$ factors as
$P \xrightarrow{f'} \bigoplus_{i' \in I'} M_{i'} \hookrightarrow
\bigoplus_{i \in I} M_i$
for some finite subset $I'$ of $I$.  By the finite case
$\bigoplus_{i' \in I'} M_{i'}$ is Mittag-Leffler and hence there exists
a finitely presented  module $Q$ and a map $g: P \to Q$ such that $g$
and $f'$ dominate each other.  Then also $g$ and $f$ dominate each other.
\end{proof}


\begin{lemma}
\label{lemma-split-ML-henselian}
Let $R$ be a henselian local ring.
Any countably generated Mittag-Leffler module over $R$ is a direct
sum of finitely presented $R$-modules.
\end{lemma}

\begin{proof}
Let $M$ be a countably generated and Mittag-Leffler $R$-module.
We claim that for any element $x \in M$ there exists a direct
sum decomposition $M = N \oplus K$ with $x \in N$, the module
$N$ finitely presented, and $K$ Mittag-Leffler.

\medskip\noindent
Suppose the claim is true. Choose generators $x_1, x_2, x_3, \ldots$
of $M$. By the claim we can inductively find direct sum decompositions
$$
M = N_1 \oplus N_2 \oplus \ldots \oplus N_n \oplus K_n
$$
with $N_i$ finitely presented,
$x_1, \ldots, x_n \in N_1 \oplus \ldots \oplus N_n$, and $K_n$ Mittag-Leffler.
Repeating ad infinitum we see that $M = \bigoplus N_i$.

\medskip\noindent
We still have to prove the claim. Let $x \in M$. By
Lemma \ref{lemma-ML-countable}
there exists an endomorphism $\alpha : M \to M$
such that $\alpha$ factors through a finitely presented module, and
$\alpha (x) = x$. Say $\alpha$ factors as
$$
\xymatrix{
M \ar[r]^\pi & P \ar[r]^i & M
}
$$
Set $a = \pi \circ \alpha \circ i : P \to P$, so
$i \circ a \circ \pi = \alpha^3$. By
Lemma \href{https://stacks.math.columbia.edu/tag/05BT}{lemma charpoly module (Tag 05BT)}
there exists a monic polynomial $P \in R[T]$ such that $P(a) = 0$.
Note that this implies formally that $\alpha^2 P(\alpha) = 0$.
Hence we may think of $M$ as a module over $R[T]/(T^2P)$.
Assume that $x \not = 0$. Then $\alpha(x) = x$ implies that
$0 = \alpha^2P(\alpha)x = P(1)x$ hence $P(1) = 0$ in $R/I$ where
$I = \{r \in R \mid rx = 0\}$ is the annihilator of $x$.
As $x \not = 0$ we see $I \subset \mathfrak m_R$, hence
$1$ is a root of $\overline{P} = P \bmod \mathfrak m_R \in R/\mathfrak m_R[T]$.
As $R$ is henselian we can find a factorization
$$
T^2P = (T^2 Q_1) Q_2
$$
for some $Q_1, Q_2 \in R[T]$ with
$Q_2 = (T - 1)^e \bmod \mathfrak m_R R[T]$ and
$Q_1(1) \not = 0 \bmod \mathfrak m_R$, see
Lemma \href{https://stacks.math.columbia.edu/tag/04GG}{lemma characterize henselian (Tag 04GG)}.
Let $N = \Im(\alpha^2Q_1(\alpha) : M \to M)$ and
$K = \Im(Q_2(\alpha) : M \to M)$. As $T^2Q_1$ and
$Q_2$ generate the unit ideal of $R[T]$ we get a direct sum
decomposition $M = N \oplus K$. Moreover, $Q_2$ acts as zero on $N$ and
$T^2Q_1$ acts as zero on $K$. Note that $N$ is a quotient of $P$
hence is finitely generated. Also $x \in N$ because
$\alpha^2Q_1(\alpha)x = Q_1(1)x$ and $Q_1(1)$ is a unit in $R$. By
Lemma \ref{lemma-direct-sum-ML}
the modules $N$ and $K$ are Mittag-Leffler. Finally, the finitely generated
module $N$ is finitely presented as a finitely generated Mittag-Leffler
module is finitely presented, see
Example \href{https://stacks.math.columbia.edu/tag/059R}{example ML (Tag 059R)} part (1).
\end{proof}
\end{document}
